%==============================================================================
%  Appendix D  Extension: the Hierarchical Block Model              [owner: A1]
%  Source: v8 app:hbm (v8:1005-1117).
%  asm:hbm and tab:assumption_chain live HERE (they were in sec:problem between the
%  v9 refine pass and the CFBM audit; moved back because no theorem cites asm:hbm).
%  There must be exactly ONE \label{asm:hbm} in the document -- a second copy left
%  behind in sec:problem would silently double-number every later assumption.
%  Retargeted: \Cref{lem:A-gap} -> \Cref{lem:gap} (twice).
%==============================================================================
% [Reviewer Note: this appendix used to open "asm:cfbm is replaced by asm:hbm",
%  which overstated what is at risk.  Since asm:cfbm was restated in floor form
%  (eq:cfbm_floor), the HBM no longer competes with it: the HBM keeps the cross-clan
%  block constant and only constrains the diagonal blocks, which eq:cfbm_floor
%  leaves free apart from a floor.  So the HBM SATISFIES asm:cfbm rather than
%  replacing it -- the Fiedler vector stays exactly piecewise constant, thm:main-sim
%  applies unchanged, and the only quantity that moves is the gap in lem:gap.  The
%  appendix is therefore a corollary, not an extension.  This also settles register
%  item R1/04: with the floor stated directly, structural collapse has instances.]
\section{The Hierarchical Block Model as a sufficient condition}\label{app:hbm}
The floor $\Sin^{\min}$ of \eqref{eq:cfbm_floor} is an abstract property of the
matrix. This appendix connects it to the tree: under geometric per-edge decay the
floor is pinned by the depth of the deeper clan, which is what makes the
gap --- and therefore recoverability --- a function of tree shape rather than of a
free scalar. The hypothesis that does it is the following.

\begin{assumption}[Hierarchical Block Model with geometric decay]\label{asm:hbm}
The HBM instantiates the floor of \eqref{eq:cfbm_floor} by letting intra-clan
similarities decay geometrically with
path depth, parameterized by the maximum tree depth $D_{\max}$, a per-edge decay
factor $\alpha\in(0,1]$ (under GTR with unit edges $\alpha=\exp(4\,\mathrm{tr}(Q)\,\Delta\tau)$),
and a top-level baseline $\Sin$. We say $S$ satisfies the HBM if the weakest
intra-clan similarity is bounded below by
$\Sin^{\min}\ge\Sin\,\alpha^{D_{\max}}$. The depth $D_{\max}$ is governed by tree
shape: under Yule or Kingman models it scales as $\mathcal O(\log m)$, while
degenerate caterpillar geometries reach $\mathcal O(m)$
\cite{steel2001properties, blum2006which}.
\end{assumption}

\Cref{asm:hbm} is \emph{added} to \Cref{asm:gtr,asm:clock,asm:cfbm}, none
of which is withdrawn. It keeps the constant cross-clan boundary and constrains only
the diagonal blocks --- exactly the part that \eqref{eq:cfbm_floor} leaves free apart
from the floor. It is therefore a sufficient condition for \Cref{asm:cfbm} rather
than a substitute for it: under the HBM the Fiedler vector is still exactly
piecewise constant (\Cref{lem:A-decomp}), \Cref{thm:main-sim} applies verbatim, and
the only quantity that moves is the spectral gap. No result in the main text takes
\Cref{asm:hbm} as a hypothesis. \Cref{tab:assumption_chain} places it in the chain
that runs from the generative model to the block structure of $S$; the rest of the
appendix draws the two consequences, a geometric floor on the weakest intra-clan
similarity and the spectral-gap bound that floor induces.

\begin{table}[htbp]
\centering\footnotesize\renewcommand{\arraystretch}{1.2}
\begin{tabular}{@{}p{0.26\textwidth} p{0.66\textwidth}@{}}
\toprule
\textbf{Generative assumption} & \textbf{Consequence for $S$} \\
\midrule
GTR alone & $S_{ij}=\exp(4\,\mathrm{tr}(Q)\,\tau_{ij})$; diagonal (intra-clan)
blocks dominate off-diagonal blocks by MRCA-depth monotonicity. \\
\addlinespace
GTR + molecular clock & Localised CFBM: cross-clan similarities decay uniformly to
an exact baseline $\bze$. \\
\addlinespace
GTR + clock + intra-clan floor & Full CFBM: within-clan similarities clear the
baseline by $\Sin^{\min}-\bze>0$, their internal structure left free; applies to
symmetric and asymmetric trees alike. \\
\addlinespace
GTR + clock + path-depth variance & HBM: intra-clan similarities decay
geometrically as $\alpha^{\mathrm{depth}}$, which supplies the floor
$\Sin^{\min}\ge\Sin\,\alpha^{D_{\max}}$ at depth $D_{\max}$. \\
\bottomrule
\end{tabular}
\caption{From the generative model to the block structure of $S$.}
\label{tab:assumption_chain}
\end{table}

\begin{proposition}[Geometric floor]\label{prop:D-floor}
Under \Cref{asm:hbm}, $\Sin^{\min}\ge\Sin\,\alpha^{D_{\max}}$ with
$D_{\max}=\lceil\log_2\max(a,b)\rceil$.
\end{proposition}
\begin{proof}
GTR multiplicativity along paths gives
$S_{ij}=\prod_{e\in\mathrm{path}(i,j)}\alpha_e$, so a same-clan pair at depth
$\le D_{\max}$ satisfies $S_{ij}\ge\Sin\,\alpha^{D_{\max}}$; minimising over
same-clan pairs gives the floor, with $D_{\max}$ the height of the deeper clan
subtree.
\end{proof}

\begin{proposition}[HBM spectral-gap bound]\label{prop:D-gap}
Substituting \Cref{prop:D-floor} into \Cref{lem:gap} gives
\[
   \Dlam\ge m_{\min}\bigl(\Sin\alpha^{D_{\max}}-\Sout^{\max}\bigr)-\eta\,m_{\min}\Sout^{\max}
        =m_{\min}\Sin\alpha^{D_{\max}}-(1+\eta)m_{\min}\Sout^{\max},
\]
positive iff $\alpha^{D_{\max}}>(1+\eta)\Sout^{\max}/\Sin$. As $\alpha\to1$ (or
$D_{\max}=0$) this reduces to \Cref{lem:gap} and the CFBM margin condition of
\Cref{asm:margin}.
\end{proposition}

\Cref{fig:hbm_spectral_verification} verifies the bound on simulated HBM matrices,
and \Cref{fig:fiedler_partitions} illustrates the imbalance-driven degradation of
the recovered partition.

\begin{figure}[htbp]
    \centering
    \begin{minipage}[b]{0.49\textwidth}
        \centering\includegraphics[width=\textwidth]{S_by_alpha.png}
    \end{minipage}\hfill
    \begin{minipage}[b]{0.49\textwidth}
        \centering\includegraphics[width=\textwidth]{spectral_gap_bound.png}
    \end{minipage}
    \caption{Empirical verification of \Cref{prop:D-gap} over a sweep of
    $(\eta,m)$ across decay rates $\alpha\in\{1.0,0.99,0.95,0.90\}$.
    \textbf{(Left)} HBM similarity matrices $S$ degrading as $\alpha$ decreases.
    \textbf{(Right)} slack of the gap lower bound versus $\eta$, resolving
    $\Sin^{\min}=\Sin\alpha^{D_{\max}}$ at
    $D_{\max}=\lceil\log_2\max(n_1,n_2)\rceil$; the bound stays valid
    ($\text{slack}\ge0$) across all configurations.}
    \label{fig:hbm_spectral_verification}
\end{figure}

\begin{figure}[htbp]
    \centering
    \begin{subfigure}[b]{0.48\textwidth}
        \centering\includegraphics[width=\textwidth]{Fiedler_Bipartitions/fiedler_tree_partition_eta01.png}
        \caption{$\eta=1$}\label{fig:eta01}
    \end{subfigure}\hfill
    \begin{subfigure}[b]{0.48\textwidth}
        \centering\includegraphics[width=\textwidth]{Fiedler_Bipartitions/fiedler_tree_partition_eta05.png}
        \caption{$\eta=5$}\label{fig:eta05}
    \end{subfigure}

    \vspace{0.4cm}

    \begin{subfigure}[b]{0.48\textwidth}
        \centering\includegraphics[width=\textwidth]{Fiedler_Bipartitions/fiedler_tree_partition_eta10.png}
        \caption{$\eta=10$}\label{fig:eta10}
    \end{subfigure}\hfill
    \begin{subfigure}[b]{0.48\textwidth}
        \centering\includegraphics[width=\textwidth]{Fiedler_Bipartitions/fiedler_tree_partition_eta15.png}
        \caption{$\eta=15$}\label{fig:eta15}
    \end{subfigure}
    \caption{Fiedler bipartitions of a Kingman-model tree with $500$ terminal
    nodes ($n=10{,}000$, $d=4$, JC69), across increasing imbalance $\eta$.}
    \label{fig:fiedler_partitions}
\end{figure}
